\clearpage
\subsection*{附录 G 本轮追踪与验收补充}

\begin{BSRTable}
\rowcolor{TableHead}\color{white}\bfseries 需求追踪 & \color{white}\bfseries 用例 模型 数据与验证 \\
购物车与会话 & FR-CRT-01 至 04、FR-AUTH-04 至 05 → UC18、UC23 → 图 9-1、27-1、27-6 → localStorage、sessions → AT-24、25。\\
排班与打卡 & FR-ATT-01 至 07、FR-SCH-01 至 07 → UC15、19、20 → 图 12-1、27、27-2、27-3、32-2 → users、employee\_profiles、shifts、attendance → AT-26 至 31。\\
维修协作 & FR-MNT-06 至 13 → UC12、13、21、22 → 图 11-1、24、25、27-4、27-5、32-1 → maintenance、maintenance\_events、units、stock\_events → AT-32 至 39。\\
跨角色同步 & FR-SYNC-01 至 06 → UC24 → 图 12-2、27-7 → 授权查询与持久化业务表，SSE 不单独落库 → AT-40 至 42。\\
角色化表示层 & FR-UX-01 至 04 → 各角色主用例 → 图 40 至 43 与第 5.3.4 节 → 页面数据源与权限接口 → AT-43、44。\\
\end{BSRTable}

\subsubsection*{G.1 可执行验收场景}

\begin{BSRTable}
\rowcolor{TableHead}\color{white}\bfseries 场景编号 & \color{white}\bfseries 给定 当 则 \\
AT-24 购物车 & 给定有效门店租期，加入装备后增减、移除并刷新；则选择恢复，单规格不能超过 5 或可用量，空清单不能预约；报价不产生占用，下单成功清空。\\
AT-25 退出与角色 & 给定不同标签页分别登录经理与技师，当经理退出；则经理页回登录且旧令牌无效，技师独立会话仍可读取本人任务。桌面和窄屏均能找到退出入口。\\
AT-26 角色打卡 & 给定五种非顾客角色，当各自上班再下班；则只写本人记录；顾客调用返回 403，不能通过提交其他 userId 代替本人。\\
AT-27 重复与先后 & 给定当天无记录直接下班，或已上班重复上班、已下班重复下班；则全部冲突拒绝且不多出记录。\\
AT-28 日期时区 & 给定 UTC 日期与挪威当地日期不一致、或夏令时切换日，当打卡；则工作日按 Europe/Oslo，实际时长按 UTC 差计算，不用浏览器时钟。此边界尚需专项自动化补测。\\
AT-29 班次唯一 & 给定同一人员日期已有班次，当再创建或编辑到冲突目标；则返回冲突。原日期或目标日期已有打卡时，不得编辑。\\
AT-30 日期时间 & 给定 2 月 31 日、开始晚于结束、05:59 或 23:01；则拒绝。有效 06:00 开始、23:00 结束可保存；不把门店营业窗口与排班窗口混为一谈。\\
AT-31 每日到岗 & 给定经理打开今天，店员另页上班再下班；则经理的同一人员行依次显示在岗、已下班，显示实际时间且数量一致。无班或未打卡不被自动解释为旷工。\\
AT-32 基本档案 & 给定维修技师身份，当进入维修页；则工号、姓名、岗位、联系和能力可查；缺项显示未登记，演示资质不伪装成真实审核。\\
AT-33 维修权限 & 给定工单未分配给当前技师，或店员尝试维修验收；则不返回他人工单正文或拒绝验收，不能通过隐藏按钮外的 API 绕过。\\
AT-34 提交门槛 & 给定工单在检查或维修中，当诊断、结果、正工时或四项安全检查任一缺失时提交；则拒绝并维持隔离。全部满足后只进入待验收。\\
AT-35 配件等待 & 给定维修中工单，当转待配件但原因不足 3 字；则拒绝。填写原因后可挂起，恢复先回维修，不能从待配件直接完成。\\
AT-36 验收入库 & 给定待验收且全部安全检查通过、无其他开放工单，当经理通过；则工单完成、单件可用、库存流水与审计一起出现，另一角色页面自动重查后同步。\\
AT-37 退回处理 & 给定待验收，当经理填原因退回；则回到维修中并继续隔离；技师补充处理后可以再次提交验收，原退回意见留在时间线。\\
AT-38 报废 & 给定待验收且实际无法修复，当经理确认报废；则工单与单件报废，不能再次保存为可用或被顾客预约。\\
AT-39 旧版本 & 给定两人读到版本 n，当第一人保存到 n+1 后第二人提交 n；则返回 409，不覆盖新记录，不静默丢弃用户草稿。\\
AT-40 草稿保护 & 给定一页有未保存表单，另一角色更新业务；则显示待更新但原输入保留，保存或关闭后再读最新数据。\\
AT-41 重连恢复 & 给定页面断开推送期间业务变化，当连接恢复或页面重新获得焦点；则重新读取数据库最新事实，不重放支付、打卡或维修写命令。\\
AT-42 推送内容 & 给定任一成功业务变更，当读取事件流；则仅含修订号、主题与时间，不含用户资料、令牌或业务正文；后续查询仍验证会话。\\
AT-43 NOK 语义 & 给定材料费 75、租金或押金任一金额；则明确标注 NOK；切换参考币种不改已确认金额，也不解释为人民币。\\
AT-44 页面层次 & 给定经理、维修、打卡、库存页；则各自首要任务不同且可辨认；390 像素下无整页横向溢出，按钮图标不撑高布局，退出、打卡和购物车按角色可达。\\
\end{BSRTable}

\subsubsection*{G.2 当前验证证据与限制}

本轮已完成 8 项后端集成检查：预约到归还与模拟退款、损坏复核及退款幂等、区间防超卖及价格历史、角色与调拨故障隔离、全部非顾客打卡及每日到岗、排班创建编辑与历史保护、维修安全门槛及经理验收和版本冲突、库存变更推送。浏览器检查使用隔离数据库，完成多标签页独立角色、经理每日排班、维修至验收、另一角色库存自动更新、打卡至经理到岗更新、390 像素布局及顾客购物车入口。

上述证据证明已检查路径，不表示附录所有边界或性能场景均通过。AT-28 的夏令时专项、AT-38 的报废端到端、AT-40 的草稿交互回归、AT-41 的故障注入，以及生产认证、容量、外部支付、消息重试、实时汇率和备份恢复仍需后续验证。不得把目标延迟或目标并发数写成已达到结果。

\subsubsection*{G.3 本轮尚待业务方确认的参数}

\begin{BSRTable}
\rowcolor{TableHead}\color{white}\bfseries 决策事项 & \color{white}\bfseries 当前处理与待确认内容 \\
打卡适用人员 & 按本轮要求覆盖全部非顾客角色；分包商纳入出勤记录不自动构成雇佣、薪资或纪律制度。\\
班次与考勤制度 & 当前一人同日单班、06:00–23:00、迟到早退无宽限；跨夜、多段、请假、补签及宽限分钟需组织制度决定。\\
维修费用与证据 & 当前累计工时和整数 NOK 材料费；历史人工费率、15 分钟整倍数服务端校验、照片附件和证书审核仍需完善。\\
维修验收职责 & 本轮由门店经理完成通过、退回或报废；是否需要双人复核及技师资质分级，应按 BSR 实际责任制度决定。\\
商业价格与外部服务 & 挪威 NOK 价目、押金与取消政策需正式批准；真实收退款、汇率和短信邮件接口仍为目标，不以演示记录作为已上线证据。\\
\end{BSRTable}

\subsubsection*{G.4 本轮设计参考与复现}

维修工作记录参考 ERPNext Asset Maintenance Log 的任务、负责人、期限与处理记录组织；人员打卡参考 Frappe Employee Checkin；库存变化借鉴 ERPNext Stock Ledger 的可追溯记录。BSR 的经理安全验收入库门槛按本项目责任分工另行设计，不能表述为对外部系统功能的完整复制。

\noindent\url{https://docs.frappe.io/hr/employee-checkin}\par
\noindent\url{https://docs.frappe.io/erpnext/asset-maintenance-log}\par
\noindent\url{https://docs.frappe.io/erpnext/stock-ledger}\par

\noindent 运行与检查入口：\texttt{bsr-app/README.md}；后端集成检查：\texttt{node --test test/*.test.mjs}；浏览器流程检查：\texttt{node test/preview-smoke.mjs}。界面截图来自隔离演示数据库，源代码和可编辑 LaTeX 文件均保存在本地，尚未上传 GitHub。
